\subsection{What a practitioner should take from this}

The library's answer to its own measurement is three behaviours, and the useful
part of this paper for someone fitting a process is knowing which of them they are
getting. The first is that a reported interval is a lower bound on the uncertainty
rather than the uncertainty itself. The pooled ratio is below one in each of the
twelve pooled cells, and the user who reads a 95\% interval as though it were one
gets an interval covering \ThCovSmallLow{} to \ThCovSmallHigh{} of the time across
the cells at $N = 2000$. That is a statement about the interval's width, not about
the estimate: the point estimate is unchanged, its bias is unchanged, and an
honestly wide interval around a biased estimate is still an interval around the
wrong number.

The second behaviour is refusal. In \HsixUndefined{} of the \HsixCells{} cells the
estimator is undefined, and the library reports no interval there rather than a
number extrapolated from a neighbouring regime. That refusal is the most
conservative thing it does and it is the one most likely to be mistaken for a bug,
because a fit that converges and then declines to give an interval looks broken
until one checks what the window holds. A user who hits it should check
$\beta_{\min} T$ themselves. The library does not stop them fitting such data, and
it does not warn before the fit; it reports the failure afterwards, which is
late for a user who has already spent an hour on it.

The third behaviour is the switch. Where the map measures a shortfall the fit
replaces the Hessian interval with a parametric bootstrap, and where the map is
silent it keeps the Hessian interval and says it is uncalibrated. Both are
defensible and neither is free: the bootstrap costs seconds to a couple of minutes
at the measured configurations, and it is not available at all where the map has no
measurement, which includes everything outside two types, pair decay, a spectral
radius at or below 0.95 and $N$ between 200 and 50\,000. A practitioner working in
one of those regions is back to the Hessian interval and should know it.

\subsection{What the negative results mean}

Two of the six hypotheses were not confirmed, and they fail in different ways,
which is why both are reported rather than one. H6.3's failure is a partial one:
its threshold clause held, with coverage at the lowest branching ratio exceeding
coverage at the highest by \DiscThreeT{}, and its monotonicity clause failed,
because coverage rose from $\rho = 0.3028$ to $\rho = 0.60$ before falling toward
0.95. That is the shape of a prediction that got the endpoints right and the path
wrong, and it is what one would expect from a mechanism with two competing effects:
at low branching ratio the process is nearly Poisson and the effective sample for
the decay parameters is small, at the highest branching ratio the information is
nearly singular, and in between the estimator is better behaved than either
endpoint. The hypothesis was written as a monotone decline because that is what the
near-critical intuition suggests, and the data say the intuition holds only past
the middle of the range.

H6.4's failure is a reversal and it is the more interesting of the two. The
hypothesis was that quantising timestamps to a grid coarser than the decay time
would bias the decay estimate upward, on the reasoning that ties broken at the grid
spacing look like slower decay than the ties really are. At the \DiscFourCells{}
cells with $N \ge 10\,000$ and $c = 1$, quantisation moved the relative bias of the
decay estimate down, by between \DiscFourLow{} and \DiscFourHigh{} percent, in every
one of them. The direction was wrong and the size is large enough to be a finding
in its own right: a mechanism that moves the relative bias of the decay estimate by
tens of percent, in the opposite direction from the one predicted, is a mechanism
the field does not currently model.

A reversal is more informative than a null because a null is consistent with the
effect being absent and with the design being unable to see it, while a reversal
excludes both. The sign of the error tells us something about the mechanism, and
the direction observed here suggests that what quantisation destroys is not the
timing of the decay but the multiplicity of near-simultaneous events: collapsing
events onto a grid merges them, and a merged pair looks like one event with the
energy of two, which the likelihood can explain by sharpening the decay rather than
slowing it. That explanation is a hypothesis and not a result, since nothing in this
study separates merging from timing. It is written down because it is the
explanation the direction points to and because a reader who wants to design the
study that would test it should start from here.

\subsection{What this does not establish}

The box is narrow and it is worth naming every wall. The model is two types with a
four-type satellite of eight cells, pair-specific decay, and an exponential kernel.
A power-law kernel, a shared decay rate across types, a marked or a spatial
process, or a model with covariates in the intensity are all outside what was
measured, and the map says so rather than extrapolating. Every timing in the
capability table is from one machine, so the cost figures carry no spread across
hardware. Nothing above $10^6$ events was timed and no fit reached an hour, so the
upper limit of the measured region is not located; the largest cell in the
calibration grid is $5 \times 10^4$ events, an order of magnitude below that. The
bootstrap standard errors are not part of the calibration study at all: the
separate study measured them at three cells, not across the grid, so the auto
switch's own intervals have a measured property at three points and an assumed one
everywhere else. And the study measures one estimator, at one start, with
\texttt{se\_method="hessian"}; a multi-start fit, a profile likelihood or a
different optimiser is a different estimator whose calibration was not measured
here.

Two further limits are worth stating because they are the ones a reader is most
likely to assume away. The map answers from the nearest of \HsixCells{} cells and
reports the distance, but it answers however far away that cell is, so a
\texttt{measured} status at distance two describes a regime two grid steps from the
one measured. And every cell was run at one fixed parameter set, so the map
describes those parameters at that $(N, \rho, \beta T)$, not every process with the
same summary numbers; two processes with the same spectral radius and the same
window can have different excitation structure, and the map's answer is only as
representative as the parameter set it was measured at.

\subsection{What to measure next}

Four measurements would move this from a description to a decision rule, and each
is stated as a measurement rather than a hope. The first is power at the effect
sizes the map actually reports: the study was powered to resolve a 4.5 percent
shortcut in $q$, and a replication at $R = 2000$ in the cells where $q$ is between
0.875 and 0.977 could say whether the residual shortfall there is zero or merely
small. The second is the bootstrap's own calibration across the full grid rather
than at three cells, because the auto switch is only as good as the interval it
switches to. The third is shared decay, which is a different model with $M$ decay
rates instead of $M^2$, is outside the map, and is the natural next cell population
for a practitioner who has more types than events. The fourth is whether a sandwich
estimator closes the gap where the map says it is open: equation
\eqref{eq:sandwich} gives the target, the study's stored per-replication Hessians
are not among the committed outputs, and the measurement that would settle it is a
rerun that keeps them.
